\documentclass[11pt,a4paper]{article}
\usepackage[T2A]{fontenc}
\usepackage[utf8]{inputenc}
\usepackage[russian]{babel}
\usepackage[margin=2.3cm]{geometry}
\usepackage{graphicx}
\usepackage{booktabs}
\usepackage{xcolor}
\usepackage{hyperref}
\definecolor{copper}{HTML}{C96F4A}
\hypersetup{colorlinks=true, linkcolor=copper, urlcolor=copper}
\setlength{\parskip}{4pt}
\begin{document}
\begin{center}
{\LARGE\bfseries Исследовательская программа: машина без долговременной памяти (ES1–ES5, TS1–TS3)}\\[8pt]
{\large Монография исследовательской программы · Версия 1.0.0}\\[6pt]
{\color{copper}\small ПРОГРАММА УРОВНЯ · SPINOR}\\[4pt]
{\small Исаев Исхак Хамзатович · ORCID 0009-0003-7299-0701}\\[2pt]
{\small AB Cloud Quantum Research · DOI 10.5281/zenodo.21825394}
\end{center}

\tableofcontents
\newpage

\section{Исследовательская программа: спинор-исследование}

\noindent\textit{\color{copper}Спинор-программа доводит центральный тезис «читает, не несёт» до явных моделей: единственное состояние — фаза спинора; шаг — поворот θ. Три теоремы, пять кампаний, все PASS; компромисс «ёмкость против робастности» зафиксирован как честная граница.}

\subsection{1. Цели программы}

\begin{itemize}
  \item сформулировать и доказать свойства летучего состояния: замкнутость циклов (TS1), необратимость сброса (TS2), шаг O(1) против журнала Ω(t) (TS3);
  \item измерить модели кампаниями с предрегистрированными порогами;
  \item подготовить модели для этапа P1 процессора (память V21).
\end{itemize}

\subsection{2. Модели}

\begin{table}[htbp]\centering\small
\caption{Таблица 4.2. Три модели и их параметры}
\begin{tabular}{p{0.10\textwidth} p{0.10\textwidth} p{0.13\textwidth} p{0.15\textwidth} p{0.13\textwidth} p{0.35\textwidth}}
\toprule
\textbf{Модель} & \textbf{θ} & \textbf{N = 2k} & \textbf{Бит/шаг} & \textbf{σ*} & \textbf{Назначение}
\midrule
S1 & π/2 & 4 & 2.000 & 0.135 & максимальная робастность \\
S2 & π/3 & 6 & 2.585 & 0.089 & баланс \\
S3 & π/7 & 14 & 3.807 & 0.038 & максимальная ёмкость \\
\bottomrule
\end{tabular}
\end{table}

\subsection{3. Исполненные кампании ES1–ES5}

\begin{table}[htbp]\centering\small
\caption{Таблица 4.3. Кампании спинор-исследования}
\begin{tabular}{p{0.12\textwidth} p{0.19\textwidth} p{0.52\textwidth} p{0.13\textwidth}}
\toprule
\textbf{Кампания} & \textbf{Что измерено} & \textbf{Результат} & \textbf{Вердикт}
\midrule
ES1 & циклы & периоды 4/6/14 = 2k точно; зазор = шаг & PASS \\
ES2 & сброс & I(пре;пост) = 0; реконструкция ≈ 1/N & PASS \\
ES3 & время & наклон шага −1.5e−12; журнал ×129 & PASS \\
ES4 & задача & марковская политика решаема без памяти & PASS \\
ES5 & шум & диффузия ограничена зоной тишины & PASS \\
\bottomrule
\end{tabular}
\end{table}

Эрратумы ES2a и ES5a: пороги кампаний ES2 и ES5 были заменены до принятия вердиктов; замены опубликованы в вердикт-файлах. Компромисс log₂N против 1/k — постоянная граница: ёмкость растёт, зона тишины сужается.

\subsection{4. Следующая волна}

\begin{itemize}
  \item фазы θ = π/k для k > 7: поиск точек перегиба компромисса;
  \item компиляция задач в фазовые циклы: формальная грамматика политик;
  \item диффузионный канал с корреляциями: выход за марковское предположение.
\end{itemize}

\subsection{5. Честные границы}

Не всякая задача решаема без памяти: класс применимости — задачи с марковской политикой. Реконструкция ≈ 1/N — уровень случайности, не «почти ноль»: информация стирается до предела. Журнал даёт Ω(t) при любом «умном» планировании — это нижняя граница, а не дефект реализации.

\subsection{6. Разбор кампаний}

\subsection{6.1. Пять измерений летучести}

\subsubsection{ES1 — циклы}

Первая кампания измерения моделей: периоды поворотов θ = π/2, π/3, π/7 равны 4, 6 и 14 — точно N = 2k, зазор между различимыми состояниями равен шагу 2π/N. Летучая фаза оказалась точным счётчиком: никакого дрейфа, никакой памяти.

\subsubsection{ES2 — сброс и стирание}

Сброс фазы стирает всё: взаимная информация пре/пост равна нулю, реконструкция 0.2501/0.1674/0.0717 ≈ 1/N — уровень случайности. Время после сброса определено только по модулю N. Теорема TS2 здесь получает численное тело; эрратум ES2a зафиксировал замену порога до вердикта.

\subsubsection{ES3 — реальное время}

Контрастный эксперимент: машина без памяти держит шаг O(1) (наклон −1.5e−12 мс/шаг — ноль), машина с вечным журналом растёт ×129 и срывает дедлайны. Это численное доказательство TS3 и главный аргумент архитектуры «читает, не несёт»: журнал — не надёжность, а причина отказа.

\subsubsection{ES4 — задача без памяти}

Проверка применимости: последовательность решений в цикле без хранения; на классе задач с марковской политикой качество сопоставимо с машиной с памятью. Граница честно зафиксирована: задачи с существенной историей требуют другого уровня архитектуры.

\subsubsection{ES5 — шум и зона тишины}

Диффузия фазы ограничена зоной тишины σ*: S1 (0.135 рад) устойчивее, S3 (0.038 рад) ёмче — компромисс ёмкость/робастность измерен. Эрратум ES5a: порог заменён до принятия вердикта, замена опубликована.

\subsection{7. Полная таблица кампаний}

\begin{table}[htbp]\centering\small
\caption{Полная таблица кампаний уровня}
\begin{tabular}{p{0.12\textwidth} p{0.25\textwidth} p{0.60\textwidth}}
\toprule
\textbf{Кампания} & \textbf{Название} & \textbf{Ключевой результат}
\midrule
ES1 & Фазовые циклы &  \\
ES2 & Летучесть состояния &  \\
ES3 & Реальное время &  \\
ES4 & Вычисление без хранилища &  \\
ES5 & Шумовой канал фазы &  \\
\bottomrule
\end{tabular}
\end{table}

\subsection{8. Как воспроизвести}

\begin{verbatim}
cd research/realtime_spinor_memory
python3 experiments/run_all.py        # ES1–ES5 (~2.5 мин, все PASS)
python3 make_figures.py               # плитки 4 измерений
python3 make_theorem_docs.py          # 12 документов TS1–TS3
python3 -m pytest tests/ -q           # 11 тестов
\end{verbatim}

\subsection{9. Что читать дальше}

\begin{itemize}
  \item теория: docs/theory.md (RU) и docs/theory.en.md (EN);
  \item реализация: spinor/core.py — ядро моделей, spinor/tasks.py — задачи;
  \item вердикты: data/es1…es5\_verdict.json;
  \item теоремы: theorems/TS1\_cycle, TS2\_volatility, TS3\_realtime.
\end{itemize}

\end{document}
